% -----------------------------------------------------------------------------
% One-loop QED vertex and Schwinger anomalous magnetic moment in strict SI
% -----------------------------------------------------------------------------

真空极化决定电荷随尺度的运行，电子自能决定极点质量与留数；在壳电磁顶角还包含一个不能由 Ward--Takahashi 恒等式单独固定的有限结构——Pauli 形状因子。它在零动量转移处直接给出电子反常磁矩。

\begin{figure}[htbp]
 \centering
 \begin{tikzpicture}[scale=.95,baseline=(current bounding box.center),
 fermion/.style={thick,postaction={decorate},decoration={markings,mark=at position .58 with {\arrow{Latex}}}},
 photon/.style={thick,decorate,decoration={snake,amplitude=1.2pt,segment length=6pt}}]
\coordinate (L) at (-2.2,0);
\coordinate (V1) at (-.8,0);
\coordinate (V2) at (0,0);
\coordinate (V3) at (.8,0);
\coordinate (R) at (2.2,0);
\draw[fermion] (L)--(V1)--(V2)--(V3)--(R);
\draw[photon] (V1) .. controls (-.55,1.55) and (.55,1.55) .. (V3);
\draw[photon] (V2)--(0,-1.8);
\node[above left] at (L) {$p$};
\node[above right] at (R) {$p'$};
\node[right] at (.08,-1.45) {$q=p'-p$};
\node[above] at (0,1.65) {$\ell$};
\end{tikzpicture}

 \caption{一圈 QED 顶角修正。外部电子四动量为 \(p^\mu,p'^\mu\)，外部光子注入 \(q^\mu=p'^\mu-p^\mu\)，内部光子圈动量为 \(\ell^\mu\)。所有四动量均采用 SI 物理四动量。}
 \label{fig:vertex-correction-r14}
\end{figure}

对在壳外腿，Gordon 恒等式把 Dirac 电流分成动量流与自旋磁矩两部分：
\begin{equation}
 \boxed{
 \bar u(p')\gamma^\mu u(p)
 =\frac1{2\me c}\bar u(p')
 \left[(p'+p)^\mu+\ii\sigma^{\mu\nu}q_\nu\right]u(p).}
 \label{eq:gordon-identity-r14}
\end{equation}
因此满足 Lorentz 协变性与电流守恒的在壳顶角可写成
\begin{equation}
 \boxed{
 \bar u(p')\Gamma^\mu(p',p)u(p)
 =\bar u(p')\left[
 F_1(q^2)\gamma^\mu
 +\frac{\ii\sigma^{\mu\nu}q_\nu}{2\me c}F_2(q^2)
 \right]u(p).}
 \label{eq:form-factor-definition-r14}
\end{equation}
其中 \(F_1(0)=1\) 由物理电荷归一化固定，\(F_2(0)\) 给出反常磁矩。

Feynman 规范中的一圈约化顶角修正从
\begin{equation}
 \Lambda_{\rm 1L}^\mu(p',p)
 =-\ii g_{\rm em}^2
 \int_{\ell}^{(d)}
 \frac{N^\mu(\ell)}{D_\gamma D'_eD_e}
 \label{eq:vertex-loop-start-r14}
\end{equation}
开始，其中三个传播子分母和自旋子分子分别为
\begin{equation}
\boxed{
\begin{aligned}
D_\gamma&=\ell^2+\ii0^+,\\
D'_e&=(p'-\ell)^2-\me^2c^2+\ii0^+,\\
D_e&=(p-\ell)^2-\me^2c^2+\ii0^+.
\end{aligned}}
\label{eq:vertex-propagator-denominators-dp75}
\end{equation}
\begin{equation}
\boxed{
N^\mu(\ell)=\gamma^\alpha
(\slashed p'-\slashed\ell+\me c)\gamma^\mu
(\slashed p-\slashed\ell+\me c)\gamma_\alpha.}
\label{eq:vertex-loop-numerator-dp75}
\end{equation}
三个分母依次对应内部光子以及外部光子顶角两侧的电子传播子；$N^\mu$ 是内部电子传播子的 Dirac 分子与三个 QED 顶角收缩后留下的自旋子空间矩阵。对\eqnrefs{eq:feynman-three-denominator-r14} 取
$A=D'_e$、$B=D_e$、$C=D_\gamma$，并利用在壳条件
$p^2=p'^2=\me^2c^2$。定义平移后的圈动量与参数化质量函数
\begin{equation}
\boxed{
L^\mu=\ell^\mu-xp'^\mu-yp^\mu,\qquad
\Delta_V(x,y,z;q^2)=(1-z)^2\me^2c^2-xyq^2,\qquad x+y+z=1,}
\label{eq:vertex-shift-delta-dp75}
\end{equation}
便有
\begin{equation}
\boxed{
\Lambda_{\rm 1L}^\mu
=-2\ii g_{\rm em}^2
\int_0^1\dd x\dd y\dd z\,
\delta(x+y+z-1)
\int_L^{(d)}
\frac{N^\mu(L;x,y,z)}{(L^2-\Delta_V+\ii0^+)^3}.}
\label{eq:vertex-parameterized-mother-dp75}
\end{equation}
\eqnrefs{eq:vertex-parameterized-mother-dp75} 是后续顶角计算的母积分；$\Delta_V$ 具有动量平方量纲，并把全部外部运动学集中在 Feynman 参数上。进一步的 Clifford 代数把该积分投影到\eqnrefs{eq:form-factor-definition-r14} 的两个独立 Lorentz 结构。对 Pauli 结构而言，参数化后的分子最终只留下
\begin{equation}
 \boxed{N_2(x,y,z)=4z(1-z)\me^2c^2,\qquad x+y+z=1.}
 \label{eq:vertex-pauli-numerator-r68}
\end{equation}


取零动量转移后，一圈 Pauli 形状因子得到 Schwinger 结果
\begin{equation}
 \boxed{
 F_2^{\rm 1L}(0)=\frac{\alpha}{2\pi}.}
 \label{eq:schwinger-F2-r14}
\end{equation}
Pauli 形状因子通过旋磁因子与实验磁矩相连。旋磁因子的物理定义为：对自旋 $1/2$ 粒子，$\mathbf S_{1/2}=\hbar\boldsymbol\sigma_{\rm P}/2$；把均匀弱磁场中的自旋能量写成
\begin{equation*}
 H_B=-\frac{g_{\rm mag}\qe}{2\me}\,\mathbf S_{1/2}\cdot\mathbf B,
\end{equation*}
其中无量纲数 $g_{\rm mag}$ 称为旋磁因子。树级 Dirac 方程给 $g_{\rm mag}=2$，而量子顶角修正会改变这个系数。这个 $g_{\rm mag}$ 与电磁顶角耦合常数 $g_e=\qe/\sqrt{\varepsilon_0\hbar c}$ 是不同物理量。

物理电荷归一化给出 \(F_1(0)=1\)。把偏离 Dirac 值 $2$ 的部分写成无量纲反常磁矩系数，则
\begin{equation}
 \boxed{
 a_\ee^{\rm anom}
 \equiv\frac{g_{\rm mag}-2}{2}
 =F_2(0),
 \qquad
 g_{\rm mag}=2[1+F_2(0)].}
 \label{eq:ae-form-factor-r68}
\end{equation}
结合\eqnrefs{eq:schwinger-F2-r14}，一圈阶有 \(a_\ee^{\rm anom}=\alpha/(2\pi)\)。普通一圈截断由 \(\alpha/\pi\ll1\) 控制；更高精度需要系统加入二圈及以上顶角、自能和真空极化贡献。

保留类空动量转移 \(Q^2\equiv-q^2>0\) 时，一圈结果可化为
\begin{equation}
 \boxed{
 \frac{F_2^{\rm 1L}(-Q^2)}{F_2^{\rm 1L}(0)}
 =\int_0^1\dd u_F\,
 \frac1{1+[Q/(\me c)]^2u_F(1-u_F)}.}
 \label{eq:F2-spacelike-ratio-r14}
\end{equation}
唯一的运动学比值是 \(Q/(\me c)\)：低动量极限恢复静态磁矩，高动量时电子有限 Compton 尺度使 Pauli 形状因子逐渐受抑。

\begin{figure}[htbp]
 \centering
 \includegraphics[width=.76\linewidth]{qed_pauli_form_factor_r14.png}
 \caption{一圈 QED Pauli 形状因子的类空动量依赖。横轴为 \(Q/(\me c)\)，纵轴以 \(F_2(0)=\alpha/(2\pi)\) 归一化；曲线由\eqnrefs{eq:F2-spacelike-ratio-r14} 数值积分得到。}
 \label{fig:qed-pauli-form-factor-r14}
\end{figure}
